\chapter{Elasticity}
\label{chap:elasticity}
%\section{Theory}
The Elasticity module computes the mechanical equilibrium equation 
\begin{equation}
\frac{ \partial \sigma_{ij}}{\partial x_j} = f_i
\end{equation}
where $\mathbf{f}$ is the total body force and  $\sigma$ is the stress tensor given by 
\begin{equation}
\sigma_{ij} =C_{ijlk} \epsilon_{lk}
\end{equation}
$C$ is the fourth-rank stiffness tensor and $\epsilon$ is the strain tensor expressed in terms of displacements $u$ by
\begin{equation}
\epsilon_{lk} =\frac{1}{2} \left( \frac{\partial u_l}{\partial x_k} + \frac{\partial u_k}{\partial x_l}\right)
\end{equation}

\section{Options}

The option \emph{shape\_itarations} sets the number of iterations for the strain calculation (1 means no deformation and is the default). The option \emph{shape\_error} is max allowed error for the deformation (default = 1e-2).


\section{Models}
Scheme:
\begin{verbatim}
Physical_model Name
{
   Type = 
   Options
}
\end{verbatim}
\noindent Bulk: apply for physical regions wit dim = N\\
Inteface: apply for physical regions with dim = N-1\\
N: dimension of the mesh.
\subsection{Stiffness/Isotropic} 
Kind: bulk \\
Name: \emph{stiffness}\\
Type: \emph{isotropic}\\
Options:  \emph{young = 129} (Young modulus), \emph{poisson = 0.349} (Poisson ration) 

\subsection{Stiffness/Anisotropic (default)}
Kind: bulk \\
Name: \emph{stiffness}\\
Type: \emph{anisotropic} 

\subsection{Body force/constant}
Kind: bulk \\
Name: \emph{body\_force}\\
Type: \emph{constant}\\
Options: \emph{force = (1,0,0)}\\

\subsection{Body force/Lattice mismatch}
Kind: bulk \\
Name: \emph{body\_force}\\
Type: \emph{lattice\_mismatch}\\
Options: material options with the same format of the device section.

\subsection{Body force/Converse}

Kind: bulk \\
Name: \emph{body\_force}\\
Type: \emph{converse} \\
Options: \emph{poisson\_simulation}\\
Description: includes the converse piezo body force.

\subsection{Boundary/Clamp}
Kind: interface \\
Name: \emph{boundary}\\
Type: \emph{clamp} \\
Options: none \\
Description: fixes the node to have zero displacements.

\subsection{Boundary/Substrate}
Kind: interface \\
Name: \emph{boundary}\\
Type: \emph{substrate} \\
Options: none \\
Description: fixes nodes to move only along the normal to the surface they belong to.

\subsection{Boundary/Surface force}
Kind: interface \\
Name: \emph{boundary}\\
Type: \emph{surface\_force} \\
Options: \emph{force = (0,0,1)} \\
Description: applies a surface force

\section{Output}

All data are plotted on nodes. Keywords: Displacement,Strain,Stress,ForceSource,StressSource,StrainSource.




